% MAT670 -- Interlude G (placed after Interlude F): packings in the laboratory.
% Entirely new material (not in the handwritten notes). Labels prefixed PHYS:.
% Numerical claims checked with python3 (cell areas, Coxeter's count, Bragg spacings, Brillouin-zone
% distances, the free-volume model, Letokhov's radius, the n=6 sticky-cluster partition functions).

\newcommand{\PHYStag}[1]{\textsf{\scriptsize[#1]}}
\tikzset{
  PHYSth/.style={fill=black,draw=black},
  PHYSball/.style={fill=addcol!25,draw=addcol,fill opacity=0.45},
}
\newcommand{\PHYSmth}[1]{\fill[black] #1 ++(-2.2pt,-2.2pt) rectangle ++(4.4pt,4.4pt);}
\newcommand{\PHYSmsim}[1]{\filldraw[draw=black,fill=white,thick] #1 circle (2.3pt);}
\newcommand{\PHYSmexp}[1]{\fill[flagcol] #1 ++(0pt,3.2pt) -- ++(-2.9pt,-5.2pt) -- ++(5.8pt,0pt) -- cycle;}
\newcommand{\PHYSmcon}[1]{\fill[addcol] #1 ++(0pt,3pt) -- ++(3pt,-3pt) -- ++(-3pt,-3pt) -- ++(-3pt,3pt) -- cycle;}

\chapter[Interlude G: Packings in the lab]{Interlude G: Packings in the laboratory}
\label{PHYS:ch}

\begin{addedblock}
This interlude is new: it is not in the handwritten notes. It follows the aside at the end of
Lecture~6 (``ellipsoid packings; polygon packings; honeycomb and bubble problems'') and the physical
remarks scattered through Lectures~7--11 out of the mathematics and into the laboratory: squashed peas,
foams, colloids, opals, lasers made of sand, sticky spheres and dice. To keep it readable it is set in
black; all of it should be read as added material.
\end{addedblock}

Everything in these notes so far has been about ideal balls: exactly round, exactly hard, exactly
equal. Physicists have been squeezing, shaking and photographing real ones for three centuries, and
many of the objects of the course were met on a laboratory bench before they were met in a proof. The
Voronoi cell was a squashed pea, the face-centred cubic lattice was an opal, and the contact graph was
a pattern of paint marks on ball bearings.

In physics one has to be careful about the status of a statement, more careful than in the rest of the
course. Throughout, claims are tagged \PHYStag{theorem} (proved), \PHYStag{simulation} (seen on a
computer), \PHYStag{experiment} (measured) or \PHYStag{model} (a calculation in an idealized model).
Most of the famous numbers of the subject (the random close packing fraction $0.64$, the freezing
point $0.494$ of hard spheres, the jamming exponents) are of the second and third kind; Figure~\ref{PHYS:fig:strip} sorts them. Interlude~E
discussed the statistical mechanics behind several of them (hard spheres, random close packing, Maxwell
counting, frustration, Edwards' entropy), and we refer to it rather than repeat it.

%======================================================================
\section{Squashed peas, lead shot and soap}
\label{PHYS:sec:squash}

\subsection*{The first Hales}
Lecture~1 introduced the Voronoi cell (Definition~\ref{L1:def:voronoi}) as a bookkeeping device: it
assigns to each ball the part of space that is closer to it than to any other ball. There is a much
more physical way to see Voronoi cells. Take soft balls, and press them together until there is no
space left. If every ball keeps its centre and grows at the same rate, each one ends up filling exactly
its Voronoi cell.

The first famous squashing experiment was made by a Hales, although not the one who proved the
Kepler conjecture. Stephen Hales, a clergyman and physiologist, reported in his \emph{Vegetable
Staticks} (1727):
\begin{quote}\small
I compressed several fresh parcels of Pease in the same Pot, with a force equal to 1600, 800, and 400
pounds; \dots\ being thereby formed into pretty regular Dodecahedrons.
\end{quote}
(The peas swelled with water inside an iron pot, under a weighted lid.) What could ``pretty regular
dodecahedrons'' have been? Not regular pentagonal dodecahedra, as a short argument shows.

\begin{proposition}\label{PHYS:prop:dodec}
Congruent regular dodecahedra do not tile $\R^3$.
\end{proposition}

\begin{proof}
The dihedral angle of the regular dodecahedron is $\delta=\arccos(-1/\sqrt5)\approx116.57^\circ$.
Suppose the dodecahedra tile, and let $e$ be an edge of one tile $T$. The edges and vertices of the
other tiles meet the open segment $e$ in finitely many points or in subsegments of $e$, so we may pick
a point $p\in e$ that is not a vertex of any tile and at which no edge of another tile crosses $e$
transversally. Let $T'\ne T$ be a tile containing $p$. If $p$ lies in the relative interior of a face
of $T'$, that face must contain the direction of $e$: otherwise the face plane would cross $e$ at $p$
and the half-space bounded by it would overlap the interior of $T$ near $p$. If $p$ lies on an edge of
$T'$, then that edge runs along $e$. So in a small disc perpendicular to $e$ at $p$, the tiles cut out
sectors of angle $\delta$ (wedges along $e$) and $\pi$ (half-planes), with disjoint interiors, filling
the disc. Hence $k\delta+m\pi=2\pi$ for integers $k\ge1$, $m\ge0$. Since $\pi/2<\delta<\pi$, the case
$m=1$ is impossible, and so is $m=2$. For $m=0$ we would need $k\delta=2\pi$, but
$3\delta\approx349.7^\circ$ and $4\delta\approx466.3^\circ$.
\end{proof}

So three regular dodecahedra around an edge leave a gap of about $10.3^\circ$. This is the
dodecahedral counterpart of the $7.36^\circ$ gap left by five regular tetrahedra around an edge
(Section~\ref{TET:sec}), and the two facts are dual to each other: a regular dodecahedron is the
Voronoi cell of a ball surrounded by a regular icosahedron of twelve neighbours, whose Delaunay cells
would be the twenty slightly flattened tetrahedra of Interlude~E. By the theorem of Hales and
McLaughlin quoted there, no Voronoi cell of a unit ball packing is smaller than this dodecahedron, and
by Proposition~\ref{PHYS:prop:dodec} not all of them can be equal to it. The peas must have been
something else.

\subsection*{Kepler's pomegranate}
A century earlier, Kepler had already given the right answer. In \emph{Strena seu de nive sexangula}
(1611), the little book on snowflakes that also contains his packing conjecture, he describes the
seeds of a pomegranate, squeezed together as the fruit grows, and the rhombic shape they take: the
rhombic dodecahedron. Here is the computation behind it.

\begin{proposition}[The Voronoi cell of FCC]\label{PHYS:prop:rd}
Let $D_3=\{x\in\Z^3:x_1+x_2+x_3\ \text{even}\}$ be the FCC lattice, with minimal distance $\sqrt2$. Its
Voronoi cell at $0$ is
\[
  V=\{x\in\R^3:\ \abs{x_i}+\abs{x_j}\le1\ \text{for all }i\ne j\}.
\]
This is the rhombic dodecahedron: the cube $[-\frac12,\frac12]^3$ with a square pyramid of height
$\frac12$ erected on each face. It has $12$ rhombic faces, one for each minimal vector $\pm e_i\pm e_j$, $6$
vertices of degree four at $(\pm1,0,0),(0,\pm1,0),(0,0,\pm1)$, and $8$ vertices of degree three at
$(\pm\frac12,\pm\frac12,\pm\frac12)$. Its volume is $2=\covol(D_3)$.
\end{proposition}

\begin{proof}
The Voronoi cell is $\{x:\ip xv\le\frac12\abs v^2\ \forall v\in D_3\setminus0\}$. For the $12$ minimal
vectors $v=\pm e_i\pm e_j$ the condition reads $\pm x_i\pm x_j\le1$, so $V$ contains the Voronoi cell. For
the converse we show that every other $v\in D_3$ gives a redundant condition. Let $x\in V$ and put
$a_i=\abs{x_i}$; then $a_i+a_j\le1$, so each $a_i\le1$ and at most one $a_i$ exceeds $\frac12$. With
$m_i=\abs{v_i}\in\Z_{\ge0}$,
\[
  \tfrac12\abs v^2-\ip xv\ \ge\ \sum_i\Big(\tfrac12m_i^2-a_im_i\Big).
\]
A term with $m_i=0$ is $0$, and a term with $m_i\ge2$ is at least $\frac12m_i(m_i-2)\ge0$. A term with
$m_i=1$ is $\frac12-a_i$, which is negative for at most one index $i_0$. If that happens, $v_{i_0}$ is odd,
and since $v_1+v_2+v_3$ is even, some other coordinate $v_j$ is odd too. If $m_j=1$, the two terms add up
to $1-a_{i_0}-a_j\ge0$; if $m_j\ge3$, the $j$-term is at least $\frac32$ and the $i_0$-term at least
$-\frac12$. Either way the sum is $\ge0$. So $V$ is the Voronoi cell. (For every simply laced root lattice
the same holds, with the roots in place of the $\pm e_i\pm e_j$; see Lemma~\ref{DYN:lem:relevant} of
Interlude~H.)

On the cube $\max\abs{x_i}\le\frac12$ all the inequalities hold. Over the face $x_3=\frac12$ of the cube,
the inequalities $\abs{x_1}+x_3\le1$ and $\abs{x_2}+x_3\le1$ cut out the pyramid with apex $(0,0,1)$, and
similarly for the other faces. The volume is $1+6\cdot\frac13\cdot1\cdot\frac12=2$. Each face
$\{\pm x_i\pm x_j=1\}\cap V$ is a rhombus with two degree-three and two degree-four vertices, for instance
$(1,0,0),(\frac12,\frac12,\pm\frac12),(0,1,0)$ on the face $x_1+x_2=1$.
\end{proof}

The two kinds of vertices are the two kinds of holes of FCC (Interlude~H, Section~\ref{DYN:sec:fcc}): the
degree-four vertices, at distance $1$ from the centre, are octahedral holes, and the degree-three ones,
at distance $\frac{\sqrt3}{2}$, are tetrahedral holes. The rhombic faces have diagonals $\sqrt2$ and $1$.
In the hexagonal close packing (HCP) the layers above and below each ball are mirror images of each
other, and the cell becomes the \emph{trapezo-rhombic dodecahedron}: the rhombic dodecahedron cut across a
three-fold axis, with one half turned by $60^\circ$ (six rhombi and six trapezoids, the same volume). Both
cells have twelve faces, one for each of the twelve kissing neighbours. Whatever Stephen Hales saw, the
experiment below suggests that most of his peas were less regular than either.

\begin{figure}[ht]
\centering
\begin{minipage}{0.47\textwidth}
\centering
\begin{tikzpicture}[x={(1.55cm,0cm)},y={(0.62cm,0.42cm)},z={(0cm,1.55cm)},line join=round]
  \foreach \a in {0.5,-0.5} \foreach \b in {0.5,-0.5} {
    \draw[dashed,gray] (-0.5,\a,\b)--(0.5,\a,\b);
    \draw[dashed,gray] (\a,-0.5,\b)--(\a,0.5,\b);
    \draw[dashed,gray] (\a,\b,-0.5)--(\a,\b,0.5);}
  \foreach \h in {0.5,-0.5} \foreach \a in {0.5,-0.5} \foreach \b in {0.5,-0.5} {
    \draw[thick] ({2*\h},0,0)--(\h,\a,\b);
    \draw[thick] (0,{2*\h},0)--(\a,\h,\b);
    \draw[thick] (0,0,{2*\h})--(\a,\b,\h);}
  \fill[flagcol] (0,0,1) circle (2.2pt) node[above,font=\scriptsize,black] {$(0,0,1)$: octahedral hole};
  \fill[addcol] (0.5,-0.5,0.5) circle (2.2pt) node[right,font=\scriptsize,black,xshift=2pt] {$\tfrac12(1,-1,1)$: tetrahedral hole};
  \fill (0,0,0) circle (1.2pt);
\end{tikzpicture}
\end{minipage}\hfill
\begin{minipage}{0.5\textwidth}
\centering\small
\renewcommand{\arraystretch}{1.15}
\begin{tabular}{lcc}
\toprule
cell & faces & $A/V^{2/3}$\\
\midrule
ball (does not tile) & -- & $4.836$\\
rhombic dodecahedron (FCC) & $12$ & $5.345$\\
truncated octahedron (BCC) & $14$ & $5.315$\\
Kelvin (curved faces) & $14$ & $5.306$\\
Weaire--Phelan (average) & $13.5$ & $5.288$\\
cube & $6$ & $6$\\
\bottomrule
\end{tabular}
\end{minipage}
\caption{Left: the Voronoi cell of FCC is a cube (dashed) with a pyramid on each face. Right: surface area
per cell for cells of unit volume. The flat polyhedra were computed for this table; the Kelvin and
Weaire--Phelan values are numerical minimizations (Surface Evolver), as quoted by Weaire and
Phelan.}
\label{PHYS:fig:rd}
\end{figure}

\subsection*{Lead shot: what disorder does}
Hales' peas presumably went into the pot anyhow. In 1939 the botanist J.~W.~Marvin repeated the experiment with
lead shot, compressed in a cylinder \PHYStag{experiment}. Shot that had been carefully stacked in a
close-packed arrangement became rhombic dodecahedra. Shot that had simply been poured in became
irregular polyhedra with on average $14.17$ faces, not twelve. So disorder costs faces: a randomly packed ball has
more Voronoi neighbours than a crystalline one, even though it has fewer \emph{touching} neighbours.

Bernal turned this into a programme around 1960. To model the structure of liquids he poured steel
ball bearings into rubber balloons, kneaded them, ran black paint into the pile and let it set, and
then took the pile apart ball by ball, recording which pairs had been in contact by the paint marks
(Bernal, \emph{Nature} 183, 1959; Bernal and Mason, \emph{Nature} 188, 1960). Bernal and Mason counted
about $6.4$ contacts per ball \PHYStag{experiment}, close to the isostatic value $2d=6$ of
Corollary~\ref{SM:cor:maxwell}. Bernal also classified the holes of a random packing: besides the tetrahedra and octahedra of FCC there
are three larger ``canonical holes'' bounded by equilateral triangles (\emph{Proc.\ R.\ Soc.\ A} 280, 1964). The density of these packings, about $0.64$, is the
random close packing value of Interlude~E (Section~\ref{SM:sec:random}).

\subsection*{Soap: the cell that minimizes its area}
Squeeze the peas harder, or replace them by bubbles, and the cells choose their own shapes: a dry
foam minimizes its total film area subject to fixed cell volumes. Plateau (\emph{Statique
exp\'erimentale et th\'eorique des liquides}, 1873) found the rules empirically: films meet in threes, at
$120^\circ$, along curves (the Plateau borders), and the borders meet in fours at points, at the
tetrahedral angle $\arccos(-\frac13)\approx109.47^\circ$. This is the angle between two lines from the
centre of a regular tetrahedron to its vertices, the Tammes angle for $N=4$ points of
Lecture~11 (Section~\ref{L11:sec:examples}). Jean Taylor proved a century later that these are the only
singularities a soap film can have (gallery, Theorem~\ref{PHYS:thm:taylor}).

Which cell has the least area? Kelvin (1887) proposed a slightly curved truncated octahedron, the
Voronoi cell of the body-centred cubic lattice, with $8$ hexagonal and $6$ square faces. For a century
it was the best known. In 1994 Weaire and Phelan found a better foam by computer \PHYStag{simulation}:
a periodic structure with two kinds of cells of equal volume, pentagonal dodecahedra and $14$-hedra
with $12$ pentagons and $2$ hexagons, in the ratio $1:3$. Its area is about $0.3\%$ smaller
(Figure~\ref{PHYS:fig:rd}). The Weaire--Phelan foam is the Voronoi diagram, relaxed, of the A15
($\beta$-tungsten) structure, one of the \emph{Frank--Kasper} phases of metallurgy, whose Delaunay cells
are all nearly regular tetrahedra. So the best known foam is Frank's frustrated tetrahedral packing of
Interlude~E seen from the dual side. It took until 2012 to make it in the laboratory: monodisperse
bubbles tend to crystallize into Kelvin's structure next to flat walls, and Gabbrielli, Meagher, Weaire, Brakke
and Hutzler (\emph{Phil.\ Mag.\ Lett.} 92, 2012) obtained Weaire--Phelan foam in a container whose walls
were shaped to fit it \PHYStag{experiment}. Whether Weaire--Phelan is optimal is open, and so is
Kelvin's problem itself.

Matzke (\emph{Amer.\ J.\ Botany} 33, 1946) blew a foam of about a thousand equal bubbles one at a time and
examined the $600$ in its interior. He found on average $13.70$ faces per bubble and not a single Kelvin cell
\PHYStag{experiment}. Euler's formula gives a feeling for where such numbers come from.

\begin{proposition}[Coxeter's statistical honeycomb]\label{PHYS:prop:coxeter}
Let $P$ be a convex polyhedron with exactly three faces at each vertex, and let $\bar n$ be the average
number of sides of its faces. Then the number of faces is
\[
  F=\frac{12}{6-\bar n}.
\]
If every face of a foam had $\bar n=2\pi/\arccos\frac13\approx5.104$ sides, the number of regular
tetrahedra that fit around an edge, every cell would have $F\approx13.40$ faces.
\end{proposition}

\begin{proof}
With $V$ vertices and $E$ edges, $3V=2E$ (three edges at each vertex, two vertices on each edge) and $\bar
nF=2E$. Euler's formula $V-E+F=2$ becomes $F-\frac E3=2$, that is $F(1-\frac{\bar n}6)=2$. For the second
statement: by Plateau's rules every vertex of a dry foam is simple, and each face of a cell is dual to
an edge of the Delaunay complex of the cell centres. The number of sides of the face is the number of
Delaunay cells around that edge. If those cells were regular tetrahedra, the dihedral angle
$\arccos\frac13$ would have to go $\bar n$ times around the edge.
\end{proof}

Coxeter (\emph{Illinois J.\ Math.} 2, 1958) called such idealized averages ``statistical honeycombs''. The
formula is a dictionary: $\bar n=5$ gives the dodecahedron ($F=12$), Kelvin's cell has $\bar
n=\frac{8\cdot6+6\cdot4}{14}=\frac{36}7$ ($F=14$), Weaire--Phelan has $F=13.5$, Matzke's bubbles $13.70$,
and Marvin's shot $14.17$. (Voronoi cells of completely random points have on average
$2+\frac{48\pi^2}{35}\approx15.54$ faces, a theorem of Meijering, \emph{Philips Res.\ Rep.} 8, 1953.) Five tetrahedra around every edge would give
$\bar n=5$ exactly, and then all the cells would be the regular dodecahedra of
Proposition~\ref{PHYS:prop:dodec}. This is impossible in flat space, but it is exactly what happens in the
curved space of Section~\ref{SM:sec:curved}. The Voronoi cells of the $120$ vertices of the $600$-cell in
$\Sph^3$ are the $120$ regular dodecahedra of the $120$-cell, spherical dodecahedra whose dihedral angle is exactly
$120^\circ$, three around each edge. Stephen Hales' ``regular dodecahedrons'' are
what the peas would have become in a pot shaped like the $3$-sphere.

In the plane everything is settled: the hexagonal honeycomb is the least-perimeter way to divide the
plane into equal areas (Hales, gallery, Theorem~\ref{PHYS:thm:honeycomb}). Bubble rafts on water were
used by Bragg and Nye (\emph{Proc.\ R.\ Soc.\ A} 190, 1947) as models of two-dimensional crystals, complete with
dislocations and grain boundaries.

%======================================================================
\section{Colloids: atoms you can watch}
\label{PHYS:sec:colloids}

\subsection*{The hard-sphere gas, made visible}
Interlude~E (Section~\ref{SM:sec:hard}) treated hard spheres as a system with no energy at all, in which
everything is entropy. Atoms are too soft for this, but colloids, particles between about
$100$ nanometres and a few micrometres across, suspended in a liquid, come remarkably close. The standard
hard-sphere colloid is a sphere of poly(methyl methacrylate) carrying a thin brush of polymer that
stops the spheres from sticking, suspended in a liquid of nearly the same refractive index, which
removes most of the van der Waals attraction and lets light through.

Colloids are big enough to see in an optical microscope, and slow. By the Stokes--Einstein relation a sphere of radius $a$ in a liquid of viscosity $\eta$ diffuses
with $D=k_BT/(6\pi\eta a)$. For $a=1\,\mu$m in water at room temperature,
$D\approx2\times10^{-13}\,\mathrm m^2/\mathrm s$, and the time to diffuse its own radius is $a^2/D\approx5$
seconds. An atom in a liquid does the same in picoseconds. So colloids are atoms that can be filmed, and since
the only energy scale is $k_BT$, the phase behaviour depends only on the packing fraction $\varphi$.

Pusey and van Megen (\emph{Nature} 320, 1986) prepared hard-sphere suspensions at a range of
concentrations and watched them for days \PHYStag{experiment}. At low $\varphi$ the samples stayed fluid;
in a window of concentrations they separated into a fluid and a crystal, with the crystal fraction
growing across the window; above it they crystallized throughout; and above about $\varphi\approx0.58$ they
did not crystallize at all, but stayed amorphous, a colloidal \emph{glass} (Pusey and van Megen,
\emph{Phys.\ Rev.\ Lett.} 59, 1987). This is the sequence predicted by the computer for hard spheres, with coexistence between
$0.494$ and $0.545$ (Interlude~E) \PHYStag{simulation}. (The volume fractions were in fact
calibrated by identifying the onset of freezing with $0.494$, so the experiment tests the shape of the
phase diagram rather than that number; on this scale melting came out near $0.535$. Polydispersity, a few
percent in real colloids, also shifts the phase boundaries, and in simulations a size spread above about
$7\%$ suppresses crystallization altogether.) Whether hard spheres crystallize at all is still not a theorem in any dimension
$d\ge2$, as Interlude~E explained. In microgravity, on the Space Shuttle, dense samples that had stayed glassy on
earth for more than a year crystallized completely within two weeks (Zhu et al., \emph{Nature} 387,
1997), so part of the colloidal ``glass'' is gravity's doing.

\subsection*{Why order wins: a free-volume calculation}
That a system with no energy should \emph{order} when compressed seems paradoxical, since ordering
sounds like a loss of entropy. The resolution is that the relevant entropy is the room each ball has to
rattle. The following idealized model, the \emph{cell model}, makes this quantitative.

\begin{proposition}[Cell model \PHYStag{model}]\label{PHYS:prop:cell}
Suppose each of $N$ hard balls of diameter $1$ in $\R^d$ moves independently in a cage formed by its
neighbours, frozen at the sites of a lattice packing scaled from its close-packed density $\rho_{\mathrm{cp}}$
to density $\rho=N/V$, and suppose the volume of each cage is $c(a-1)^d$, where
$a=(\rho_{\mathrm{cp}}/\rho)^{1/d}$ is the nearest-neighbour distance. Then the entropy per ball is
$s=d\ln(a-1)+\ln c$ and the equation of state is
\[
  \frac{\beta P}{\rho}=\frac1{1-(\rho/\rho_{\mathrm{cp}})^{1/d}} .
\]
\end{proposition}

\begin{proof}
The partition function is $Z=(c(a-1)^d)^N$, so $S=\ln Z=Nd\ln(a-1)+N\ln c$. The volume is
$V=Na^d/\rho_{\mathrm{cp}}$, so $dV/da=Nda^{d-1}/\rho_{\mathrm{cp}}$. With $\beta P=\partial S/\partial V$ as in
Section~\ref{SM:sec:hard},
\[
  \beta P=\frac{\partial S/\partial a}{\partial V/\partial a}=\frac{Nd/(a-1)}{Nda^{d-1}/\rho_{\mathrm{cp}}}
  =\frac{\rho_{\mathrm{cp}}}{a^{d-1}(a-1)},
\]
and dividing by $\rho=\rho_{\mathrm{cp}}/a^d$ gives $\beta P/\rho=a/(a-1)=1/(1-a^{-1})$.
\end{proof}

In dimension one this is \emph{exactly} the Tonks gas of Proposition~\ref{SM:prop:tonks}: with
$\rho_{\mathrm{cp}}=1$ it reads $\beta P/\rho=1/(1-\rho)$. In one dimension the cell picture is not a
model at all, since the gaps between consecutive rods really are free lengths. In three dimensions the
cell model is crude, but not bad: for FCC at the melting density $\varphi=0.545$ it gives
$\beta P/\rho\approx10.3$, against about $11$ in simulations.

Now compare a crystal with an amorphous packing at the same density. The amorphous packing jams at
$\varphi_J\approx0.64$ instead of $\pi/\sqrt{18}\approx0.7405$, so in the same model its cages have linear size
$(\varphi_J/\varphi)^{1/3}-1$ instead of $(0.7405/\varphi)^{1/3}-1$. At $\varphi=0.55$ the ratio of the two cage sizes
is $2.01$: each crystalline ball has twice the room to move in every direction. The vibrational entropy
gained by crystallizing is therefore about $3\ln2\approx2.1\,k_B$ per ball. Against this the disordered
state has its \emph{configurational} entropy, the logarithm of the number of distinct disordered
arrangements, which wins at low density; roughly speaking, freezing is where the two balance. Order is the high-entropy state of a crowded system.

\subsection*{FCC or HCP?}
Every stacking of hexagonal layers (each layer in the hollows of the one below, the Barlow packings)
has density $\pi/\sqrt{18}$, so Hales' theorem does not choose between them. Nor does counting: with $L$
layers there are $3\cdot2^{L-1}$ stacking sequences, but each layer has $m$ balls, so the stacking entropy
per ball is $(\ln3+(L-1)\ln2)/(Lm)\to0$ as $m\to\infty$. The choice is made by the vibrational entropy,
which sees the difference between FCC and HCP only at the third shell of neighbours. Bolhuis, Frenkel,
Mau and Huse (\emph{Nature} 388, 1997) computed that FCC is favoured, by about $10^{-3}k_BT$ per particle
near melting \PHYStag{simulation}; Mau and Huse (\emph{Phys.\ Rev.~E} 59, 1999) found $0.00115(4)\,k_B$ per sphere at close
packing, and that FCC has the highest entropy of all stackings. The difference is so small that colloidal
crystals grown in the laboratory often stack \emph{randomly} (Pusey, van Megen, Bartlett, Ackerson,
Rarity and Underwood, \emph{Phys.\ Rev.\ Lett.} 63, 1989) \PHYStag{experiment}.

\paragraph{Looking inside.} Confocal microscopy locates every particle of a three-dimensional sample, in
glasses and in growing crystal nuclei (Gasser et al., \emph{Science} 292, 2001). In two dimensions, colloidal hard spheres confined to a plane
show the melting scenario that Bernard and Krauth found in simulations (Section~\ref{SM:sec:hard}): a
first-order liquid--hexatic transition followed by a continuous hexatic--solid transition (Thorneywork,
Abbott, Aarts and Dullens, \emph{Phys.\ Rev.\ Lett.} 118, 2017) \PHYStag{experiment}.

%======================================================================
\section{Opals and photonic crystals}
\label{PHYS:sec:opal}

\subsection*{A gemstone made of lattice planes}
Precious opal flashes pure colours that change as the stone is turned. In 1964 J.~V.~Sanders and his
colleagues at CSIRO looked at opal in the electron microscope and found a regular array of silica
spheres, a few hundred nanometres across, stacked in hexagonal layers (Jones, Sanders and Segnit,
\emph{Nature} 204, 1964; Sanders, \emph{Nature} 204, 1964). The colours are Bragg diffraction of light by the
layers \PHYStag{experiment}. To see which layers and which colours, we need the planes of a lattice, and
those are counted by the dual lattice $\Lat^*=\{y:\ip xy\in\Z\ \forall x\in\Lat\}$ of
Interlude~B (Theorem~\ref{NT:thm:poisson}).

\begin{proposition}[Lattice planes]\label{PHYS:prop:planes}
Let $\Lat\subset\R^3$ be a lattice and $w\in\Lat^*$ a primitive vector (not a multiple $kw'$ of
some $w'\in\Lat^*$ with $k\ge2$). Then $\Lat$ is the disjoint union of its intersections with the parallel planes
$\{x:\ip xw=k\}$, $k\in\Z$, each of which is nonempty, and these planes are $1/\abs w$ apart. Conversely, every equally spaced family of parallel planes that covers $\Lat$, each plane containing a
lattice point, arises this way, so the largest spacing between lattice planes is
$1/\min\{\abs w:w\in\Lat^*\setminus0\}$.
\end{proposition}

\begin{proof}
The values $\ip xw$, $x\in\Lat$, form a subgroup $g\Z$ of $\Z$ with $g\ge1$. Then $\ip x{w/g}\in\Z$ for all
$x\in\Lat$, so $w/g\in\Lat^*$, and primitivity forces $g=1$. So every integer $k$ is a value, every lattice
point lies on one of the planes, and the planes $\ip xw=k$ are $1/\abs w$ apart. Conversely, suppose $\Lat$ lies on the planes $\ip xu\in s\Z$, for a unit vector
$u$ and some $s>0$, and each of these planes contains a lattice point. Then $w=u/s$ has $\ip xw\in\Z$ for all
$x\in\Lat$, so $w\in\Lat^*$, and every integer is a value, so $w$ is not a multiple $kw'$ with $k\ge2$.
\end{proof}

\begin{example}[The colour of an opal]\label{PHYS:ex:opal}
For $\Lat=D_3$ the dual is $D_3^*=\Z^3\cup(\Z^3+\frac12(1,1,1))$, the body-centred cubic lattice, whose
shortest vectors are the eight vectors $\frac12(\pm1,\pm1,\pm1)$ of length $\frac{\sqrt3}2$. So the widest-spaced
planes of FCC are perpendicular to the body diagonals, at distance $2/\sqrt3$: these are the hexagonal
close-packed layers. Rescaling from minimal distance $\sqrt2$ to sphere diameter $D$, the layer spacing is
\[
  d=\frac{2}{\sqrt3}\cdot\frac{D}{\sqrt2}=D\sqrt{\tfrac23},
\]
which is also the height of a regular tetrahedron of edge $D$. The tetrahedron shows that the spacing
is the same for every Barlow stacking, so the colour does not care about stacking faults. Bragg's law for
light inside a medium of effective refractive index $n_{\mathrm{eff}}$ reflects the wavelength
\[
  \lambda=2n_{\mathrm{eff}}\,d\sin\theta,
\]
where $\theta$ is the glancing angle. At normal incidence, with $n_{\mathrm{eff}}\approx1.45$ for silica, $\lambda\approx2.37\,D$.
The visible range $400$--$700$\,nm corresponds to $D$ between about $170$ and $300$\,nm. Turning the stone
changes $\theta$, and the colour slides towards the blue.
\end{example}


\subsection*{Photonic band gaps}
In 1987 Yablonovitch (\emph{Phys.\ Rev.\ Lett.} 58, 2059), who wanted to switch off spontaneous emission, and
John (\emph{Phys.\ Rev.\ Lett.} 58, 2486), who wanted to localize light, independently proposed to do with
photons what a crystal does with electrons. In a medium whose refractive index is periodic with lattice
$\Lat$, Maxwell's equations separate into Bloch waves, and the frequencies form bands over the Brillouin
zone, the Voronoi cell of $\Lat^*$ (up to the factor $2\pi$). A \emph{complete photonic band gap} is an
interval of frequencies in which light cannot propagate in any direction.

The Brillouin zone of an opal is the Voronoi cell of BCC: the truncated octahedron, that is,
Kelvin's cell from Section~\ref{PHYS:sec:squash}. A gap opens on each face of the zone, at the frequency where
the Bragg condition for the corresponding planes holds, and for a complete gap the gaps on all faces must
overlap. With $D_3^*$ as above, the centres of the hexagonal faces (the points the physicists call $L$) are
at distance $\frac{\sqrt3}4\approx0.433$, the centres of the square faces ($X$) at $\frac12$, and the
vertices ($W$) at $\frac{\sqrt5}4\approx0.559$. The free-photon frequencies at $W$ and at $L$ differ by the factor
$\sqrt{5/3}\approx1.29$, which is hard to bridge. In fact Ho, Chan and Soukoulis (\emph{Phys.\ Rev.\ Lett.} 65,
1990) found that dielectric spheres on an FCC lattice have no complete gap in the lowest bands, because
of a band degeneracy at $W$, while spheres on the \emph{diamond} structure have one, already at an
index contrast of about $2$ \PHYStag{simulation}.

The diamond structure is $D_3^+=D_3\cup(D_3+\frac12(1,1,1))$: FCC together with one of its two classes of
tetrahedral holes (Interlude~H, Section~\ref{DYN:sec:fcc}). This is the $n=3$ case of the construction
$D_n^+$ of Lecture~6 (Section~\ref{L04b:sec:rootdense}), which is a lattice only for even $n$
(Exercise~\ref{NT:ex:D16}). As a packing of balls it is very loose: the nearest distance drops to
$\frac{\sqrt3}2$, there is one ball per unit volume, and the density is
$\frac43\pi(\frac{\sqrt3}4)^3=\frac{\pi\sqrt3}{16}\approx0.340$, less than half that of FCC. The best crystal
for light is among the worst for packing. Weevils found it first: the iridescent scales of the beetle
\emph{Lamprocyphus augustus} contain a diamond-structured photonic crystal of chitin (Galusha et al.,
\emph{Phys.\ Rev.~E} 77, 2008) \PHYStag{experiment}.

Opals can be made to work after all by inverting them. Let colloidal spheres crystallize into a synthetic
opal, fill the gaps with a material of high refractive index, and dissolve the spheres. The result, an
\emph{inverse opal} of air spheres, has a complete gap between its eighth and ninth bands once the index
contrast exceeds about $2.8$ (Busch and John, \emph{Phys.\ Rev.~E} 58, 1998) \PHYStag{simulation}.
Silicon, with index about $3.5$, qualifies, and silicon inverse opals with a complete gap near
$1.5\,\mu$m, the wavelength of optical fibres, were made by Blanco et al.\ (\emph{Nature} 405, 2000; a gap of
roughly $5\%$ of its centre frequency) and Vlasov, Bo, Sturm and Norris (\emph{Nature} 414, 2001) \PHYStag{experiment}.
For a theorem guaranteeing gaps in a model of high-contrast media, see the gallery
(Theorem~\ref{PHYS:thm:FK}).

%======================================================================
\section{From crystal to glass: random lasers and granular lasers}
\label{PHYS:sec:laser}

Take the same colloidal spheres and let them settle quickly, before they can crystallize, or add salt so
that they stick on contact. The result is a \emph{photonic glass} (Garc\'ia, Sapienza, Blanco and L\'opez,
\emph{Adv.\ Mater.} 19, 2007): a random packing of monodisperse spheres, the optical counterpart of Bernal's
ball bearings. Light no longer diffracts; it scatters many times and performs a random walk, with steps
of the order of the \emph{transport mean free path} $\ell$. Over distances much larger than $\ell$, the
energy density $W$ of the light diffuses, $\partial_tW=D\Delta W$ with $D=v\ell/3$ ($v$ the speed of light in
the medium), the continuum limit of the coin flips in the gas in a box of Lecture~7
(Section~\ref{L07:subsec:gas}).

Now soak the packing in a laser dye and pump it. Letokhov (\emph{Sov.\ Phys.\ JETP} 26, 1968) predicted
that a scattering medium with gain will lase without any mirrors, once it is large enough: the
scattering keeps the light inside long enough to be amplified. The threshold is a pretty piece of
spectral theory.

\begin{proposition}[Letokhov's critical radius \PHYStag{model}]\label{PHYS:prop:letokhov}
Let $W(r,t)$ solve $\partial_tW=D\Delta W+\frac v{\ell_g}W$ in a ball of radius $R$ in $\R^3$, with $W=0$ on the
boundary, where $D=v\ell/3$ and $\ell_g$ is the gain length (the distance over which the intensity grows
by a factor $e$). Then every solution decays to $0$ if $R<R_{\mathrm{cr}}$ and generic solutions grow
exponentially if $R>R_{\mathrm{cr}}$, where
\[
  R_{\mathrm{cr}}=\pi\sqrt{\frac{\ell\,\ell_g}3}.
\]
\end{proposition}

\begin{proof}
Expand in the Dirichlet eigenfunctions $-\Delta\phi_k=\lambda_k\phi_k$ of the ball:
\[
  W=\sum_kc_k\,e^{(v/\ell_g-D\lambda_k)t}\phi_k .
\]
Everything decays if and only if $v/\ell_g<D\lambda_1$. The first
eigenvalue is simple, with a positive eigenfunction, which is therefore invariant under rotations. For
radial $\phi$, the function $u=r\phi$ satisfies $u''=-\lambda u$ with $u(0)=u(R)=0$, so $u=\sin(\pi r/R)$ and
$\lambda_1=\pi^2/R^2$. The threshold $v/\ell_g=\frac{v\ell}3\cdot\frac{\pi^2}{R^2}$ gives the formula.
\end{proof}

The threshold volume scales like $(\ell\ell_g)^{3/2}$: stronger scattering means a smaller random laser.
And $\ell$ depends on the packing. For independent scatterers, $\ell\approx1/(n\sigma_{\mathrm{tr}})$, with
number density $n=6\varphi/(\pi D^3)$ and transport cross-section $\sigma_{\mathrm{tr}}$ of order $D^2$ for grains
much larger than the wavelength, so $\ell$ is of order $D/\varphi$ and $R_{\mathrm{cr}}\propto\varphi^{-1/2}$: compacting a
granular gain medium lowers its threshold. (Near jamming, positional correlations between the grains
reduce scattering, and the independent-scatterer estimate has to be corrected by the structure factor of
the packing.)

Lawandy, Balachandran, Gomes and Sauvain (\emph{Nature} 368, 1994) saw laser action in a dye solution
with titania particles \PHYStag{experiment}, and Cao and coworkers (\emph{Phys.\ Rev.\ Lett.} 82, 1999) saw
sharp lasing modes in zinc oxide powder, where the feedback comes from closed loops of multiply scattered
light. In a photonic glass the Mie resonances of the equal spheres can be used to select the lasing
wavelength (Gottardo et al., \emph{Nature Photonics} 2, 2008). Wiersma and Cavalieri (\emph{Nature} 414, 2001) built a random laser out
of compressed glass powder whose pores were filled with a dye-doped liquid crystal. The alignment of the
liquid crystal, and with it the strength of the scattering, depends on temperature, so the laser can be
tuned by heating it, and it switches off at the phase transition of the liquid crystal near $42^\circ$C
\PHYStag{experiment}. This was already a granular laser, tuned through the scattering of its grains.

\paragraph{Granular lasers.} Then the packing itself was made to move. Folli, Puglisi, Leuzzi and Conti
(\emph{Shaken granular lasers}, \emph{Phys.\ Rev.\ Lett.} 108, 2012) immersed about $1500$ millimetre-sized
spheres in a dye solution, put the container on a vibrating plate and pumped it with a pulsed laser
\PHYStag{experiment}. The random laser emission depends on whether and how the grains move: at strong
shaking a second emission peak appeared, attributed to transient cavities formed by a few grains. In
a sequel (Folli, Ghofraniha, Puglisi, Leuzzi and Conti, \emph{Sci.\ Rep.} 3, 2013) with $1$\,mm glass
beads in a rhodamine solution, the laser light distinguished a solid-like regime, in which the grains keep
their neighbours and the spectra are steady, from a fluidized regime in which the density fluctuations
switch the system above and below threshold, so that the spectra alternate between fluorescence and lasing.
The authors call the result a ``granular'' laser with a tunable emission spectrum: the configuration of
the beads, and whether the packing is solid-like or fluidized, is part of the optical cavity.

%======================================================================
\section{Sticky spheres}
\label{PHYS:sec:sticky}

\subsection*{From hard to sticky}
Add small polymer coils to a colloidal suspension and the spheres attract each other: when two spheres
come closer than the size of a coil, the coils are excluded from the gap between them, and the osmotic
pressure of the coils pushes the spheres together (Asakura and Oosawa, \emph{J.\ Chem.\ Phys.} 22, 1954).
The range of this \emph{depletion} attraction is the size of the coil, which can be a few percent of the
sphere diameter. In the limit of infinitely deep and narrow attraction, Baxter's \emph{sticky hard spheres}
(\emph{J.\ Chem.\ Phys.} 49, 1968), the energy of a configuration is minus the number of contacts: it is
read off the contact graph of Lectures~10--11 (Definition~\ref{L11:def:contact}), now for balls in $\R^3$.

So the ground states of $n$ sticky spheres are the packings of $n$ unit balls with the most contacts.
This \emph{contact number problem} is a cousin of the kissing problem, and it is hard. In the plane it is
solved: $n$ unit discs have at most $\lfloor3n-\sqrt{12n-3}\rfloor$ contacts (Harborth, gallery), attained by
hexagonal patches. In $\R^3$ only bounds are known (Bezdek and Reid, gallery). For small $n$ the answers
come from enumeration.

\subsection*{Rigidity and counting}
A cluster that is held together only by its contacts should be rigid. The rigidity matrix $A$ of
Interlude~E (Section~\ref{SM:sec:random}; see also Interlude~F, Section~\ref{COMP:sec:matrix}) has a row
$(\dots,-u_{ij},\dots,u_{ij},\dots)$ for each contact. Now there is no box: a free cluster can be
translated and rotated.

\begin{proposition}[Maxwell count for clusters \PHYStag{theorem}]\label{PHYS:prop:cluster}
Let $n\ge3$ unit balls in $\R^3$ have centres not all on a line. If the contact framework is
infinitesimally rigid (every infinitesimal motion preserving all contact distances to first order is a
rigid motion), then it has at least $3n-6$ contacts, and if moreover $n\ge4$, every ball has at least
three contacts.
\end{proposition}

\begin{proof}
The infinitesimal rigid motions $v_i=b+\omega\times x_i$ form a $6$-dimensional subspace of $\ker A$ when the
centres are not collinear (a motion with $b+\omega\times x_i=0$ for three non-collinear $x_i$ is zero).
Infinitesimal rigidity says this is all of $\ker A$, so $\operatorname{rank}A=3n-6$, and $A$ needs at least
that many rows. For the second claim, suppose ball $i$ has at most two contacts. Pick $w\ne0$ orthogonal
to its contact directions, and put $v_i=w$ and $v_l=0$ for $l\ne i$; this lies in $\ker A$. A rigid motion
vanishing at three non-collinear centres is zero, so if this $v$ is a rigid motion, the other $n-1\ge3$
centres lie on a line $L$, and $x_i\notin L$. Then the other balls form a straight chain along $L$, and an
end ball of the chain touches only its neighbour in the chain and possibly ball $i$. The centres other
than this end ball include $x_i$ and two points of $L$, so they are not collinear, and the same
construction at the end ball gives an element of $\ker A$ that is not a rigid motion.
\end{proof}

\noindent For clusters that are rigid but not infinitesimally rigid the second claim can fail: a ball
squeezed exactly between two balls at distance $2$ of an otherwise rigid cluster is pinned by its two
contacts.

Clusters with exactly $3n-6$ contacts that are rigid are called \emph{isostatic} (or minimally rigid),
those with more \emph{hyperstatic}, and rigid ones with fewer \emph{hypostatic}. A rigid cluster that is
not infinitesimally rigid is \emph{singular}. Arkus, Manoharan and Brenner (\emph{Phys.\ Rev.\ Lett.} 103, 2009;
\emph{SIAM J.\ Discrete Math.} 25, 2011) enumerated the candidate contact graphs (adjacency matrices with
at least $3n-6$ edges and minimum degree $3$, the two conditions suggested by Proposition~\ref{PHYS:prop:cluster}), and
solved for the distances with geometric rules. Hoy, Harwayne-Gidansky and O'Hern (\emph{Phys.\ Rev.~E} 85,
2012) went to $n=11$, and Holmes-Cerfon (\emph{SIAM Review} 58, 2016) to $n=14$ by a different method,
walking from one rigid cluster to the next along the one-parameter families obtained by breaking a
contact. Her counts:
\begin{center}
\small
\renewcommand{\arraystretch}{1.12}
\begin{tabular}{lccccccccc}
\toprule
$n$ & $6$ & $7$ & $8$ & $9$ & $10$ & $11$ & $12$ & $13$ & $14$\\
\midrule
rigid clusters & $2$ & $5$ & $13$ & $52$ & $263$ & $1659$ & $11980$ & $98529$ & $895478$\\
of which hypostatic & $0$ & $0$ & $0$ & $0$ & $1$ & $20$ & $159$ & $1308$ & $11245$\\
of which hyperstatic & $0$ & $0$ & $0$ & $0$ & $3$ & $21$ & $183$ & $1411$ & $11241$\\
most contacts & $12$ & $15$ & $18$ & $21$ & $25$ & $29$ & $33$ & $36$ & $40$\\
\bottomrule
\end{tabular}
\end{center}
These are numerical results \PHYStag{simulation}, believed complete apart possibly from a few highly
singular clusters, but not proved complete. The first singular cluster appears at $n=9$, and the first
clusters with more than $3n-6$ contacts at $n=10$: there the ground states stop being isostatic and become
fragments of the hexagonal close packing. Even the statement that $3n-6$ is the maximal number of
contacts for $6\le n\le9$ is, strictly speaking, a conjecture, conditional on the completeness of these
lists (Bezdek and Khan, arXiv:1601.00145).

\begin{figure}[ht]
\centering
\begin{tikzpicture}[scale=1.75]
  \begin{scope}
    \coordinate (O0) at (0.637,0.126); \coordinate (O1) at (-0.637,-0.126);
    \coordinate (O2) at (-0.308,0.260); \coordinate (O3) at (0.308,-0.260);
    \coordinate (O4) at (0.000,0.645); \coordinate (O5) at (0.000,-0.645);
    \foreach \i in {2,5,0,1,4,3} \filldraw[PHYSball] (O\i) circle (0.5);
    \foreach \a/\b in {0/2,0/3,0/4,0/5,1/2,1/3,1/4,1/5,2/4,2/5,3/4,3/5} \draw[thick] (O\a)--(O\b);
    \foreach \i in {0,...,5} \fill (O\i) circle (0.025);
    \node[font=\small,align=center] at (0,-1.3) {octahedron\\ $12$ contacts, $\sigma=24$};
  \end{scope}
  \begin{scope}[xshift=3.1cm]
    \coordinate (P0) at (0.661,-0.123); \coordinate (P1) at (-0.331,-0.254);
    \coordinate (P2) at (0.052,0.670); \coordinate (P3) at (0.127,0.098);
    \coordinate (P4) at (0.253,-0.856); \coordinate (P5) at (-0.762,0.466);
    \foreach \i in {3,5,4,2,0,1} \filldraw[PHYSball] (P\i) circle (0.5);
    \foreach \a/\b in {0/1,0/2,0/3,0/4,1/2,1/3,1/4,1/5,2/3,2/5,3/4,3/5} \draw[thick] (P\a)--(P\b);
    \foreach \i in {0,...,5} \fill (P\i) circle (0.025);
    \node[font=\small,align=center] at (0,-1.3) {polytetrahedron\\ $12$ contacts, $\sigma=2$};
  \end{scope}
\end{tikzpicture}
\caption{The two rigid clusters of six sticky spheres, with their contact graphs. The octahedron is a
piece of FCC (six balls around an octahedral hole). The polytetrahedron is three regular tetrahedra
glued face to face; the view is along the normal of one face of the middle tetrahedron. Both have
$3n-6=12$ contacts, so they have the same energy, but in the sticky-limit model the polytetrahedron is $24$
times as likely (Example~\ref{PHYS:ex:six}).}
\label{PHYS:fig:six}
\end{figure}

\subsection*{Six spheres: entropy chooses the less symmetric cluster}
For $n=6$ there are two ground states, both with $12$ contacts (Figure~\ref{PHYS:fig:six}). They have the same
energy, so which one forms is decided by entropy, and the entropy has three parts. Holmes-Cerfon,
Gortler and Brenner (\emph{Proc.\ Natl.\ Acad.\ Sci.\ USA} 110, 2013) showed that in the sticky limit the
partition function of a rigid, infinitesimally rigid cluster with $m$ contacts is, up to factors common
to all clusters of the same size,
\begin{equation}\label{PHYS:eq:sticky}
  z=\kappa^m\cdot\frac{\sqrt{\det I}}{\sigma\,\sqrt{\det(AA^T)}}\qquad(m=3n-6),
\end{equation}
where $\kappa$ is a ``stickiness'' parameter that depends on the depth and range of the attraction, $I$ is the
moment of inertia tensor of the centres (unit masses), $\sigma$ is the order of the rotation group of the
cluster, and $A$ is the rigidity matrix. The three factors are rotational entropy (a cluster with more
symmetry has fewer distinguishable orientations, and a more spread-out cluster has more), and vibrational
entropy: $\det(AA^T)$ is the product of the nonzero eigenvalues of the stiffness matrix $A^TA$, and
softer clusters wiggle more.

\begin{example}[$n=6$ \PHYStag{model}]\label{PHYS:ex:six}
For the octahedron, with centres $\pm\frac1{\sqrt2}e_i$, the rotation group is that of the cube,
$\sigma=24$, and $I=2\cdot\mathrm{Id}$, so $\det I=8$. The $12\times12$ matrix $AA^T$ has entries
$\ip{u_e}{u_{e'}}$ (with the signs of the shared endpoints) and $\det(AA^T)=864$. For the polytetrahedron,
with centres
\[
  (0,0,0),\ (1,0,0),\ (\tfrac12,\tfrac{\sqrt3}2,0),\ (\tfrac12,\tfrac{\sqrt3}6,\tfrac{\sqrt6}3),\
  (\tfrac12,-\tfrac{7\sqrt3}{18},\tfrac{2\sqrt6}9),\ (\tfrac43,\tfrac{4\sqrt3}9,\tfrac{2\sqrt6}9),
\]
only a half-turn survives, $\sigma=2$, and an exact computation gives $\det I=\frac{261569}{26244}\approx9.967$ and
$\det(AA^T)=\frac{261569}{972}\approx269.1$. Hence
\[
  \frac{z_{\mathrm{poly}}}{z_{\mathrm{oct}}}
  =\frac{24}2\cdot\sqrt{\frac{\det I_{\mathrm{poly}}}{\det I_{\mathrm{oct}}}\cdot\frac{\det(AA^T)_{\mathrm{oct}}}{\det(AA^T)_{\mathrm{poly}}}}
  =12\cdot\sqrt{\frac{972}{26244}\cdot\frac{864}8}=12\cdot\sqrt{\frac{108}{27}}=24 .
\]
Symmetry contributes the factor $12$, and the moments of inertia and the vibrations together exactly $2$
(about $1.12$ and $1.79$). So the model predicts that the octahedron forms with probability $\frac1{25}=4\%$.
\end{example}

Meng, Arkus, Brenner and Manoharan (\emph{Science} 327, 2010) assembled clusters of colloidal spheres with a
short-range attraction in microscopic wells and counted them \PHYStag{experiment}: the polytetrahedron
was observed more than $20$ times as often as the octahedron, so the octahedron made up only about
$4$--$5\%$ of the six-particle clusters, in good agreement with the $4\%$ (a ratio of $24$) of
Example~\ref{PHYS:ex:six}. For clusters of up to ten particles, highly symmetric clusters were strongly
suppressed by rotational entropy. So the cluster that is a piece of
the crystal loses, and the polytetrahedral one wins, as in Frank's picture of liquids (Section~\ref{SM:sec:frustration}),
though here the cause is rotational entropy rather than energy. In sticky systems the crystal comes back at
$n\ge10$, through the hyperstatic clusters, which gain a factor $\kappa$ for every extra contact. Kallus
and Holmes-Cerfon (\emph{Phys.\ Rev.~E} 95, 2017) extended \eqref{PHYS:eq:sticky} to singular clusters,
whose zero-frequency modes make the harmonic formula diverge, computed the partition functions of all
known rigid clusters with $n\le19$, and found that the hyperstatic ones, which look like fragments of a
close-packed lattice, dominate the landscape: this is how crystalline order emerges from sticky spheres,
already at $n=10$.

\paragraph{The configuration space is stratified.} The space of $n$ sticky spheres modulo rigid motions is
a union of strata, one for each contact graph $G$, of dimension $3n-6-\abs{E(G)}$ when $G$ is regular. Rigid
clusters are the $0$-dimensional strata, and breaking one contact of an isostatic cluster gives a
$1$-dimensional stratum, a path to other clusters. This is the configuration space of Lectures~7--10 (and
Interlude~D, Section~\ref{TOP:sec:conf}) cut up by contact graphs instead of by levels of the min-distance
function; see Holmes-Cerfon's review (\emph{Annu.\ Rev.\ Condens.\ Matter Phys.} 8, 2017).

%======================================================================
\section{Minimal rigidity and random jamming}
\label{PHYS:sec:jam}

\subsection*{Near the jamming point}
Interlude~E showed that a first-order jammed packing of frictionless spheres has at least $2d$ contacts
per ball on average, and that jammed packings found in simulations have exactly this many: they are
isostatic at the jamming threshold $\varphi_J\approx0.64$ (point~J of O'Hern, Silbert, Liu and Nagel). What
happens just above $\varphi_J$? Make the spheres slightly soft, with a repulsion proportional to the overlap
(harmonic spheres), and compress. Then (Figure~\ref{PHYS:fig:jam}) \PHYStag{simulation}
\[
  z-2d\ \propto\ (\varphi-\varphi_J)^{1/2},\qquad P\ \propto\ \varphi-\varphi_J ,
\]
found for a bubble model of foams by Durian (\emph{Phys.\ Rev.\ Lett.} 75, 1995) and in detail by O'Hern et
al.\ (\emph{Phys.\ Rev.~E} 68, 2003). In the laboratory, Majmudar, Sperl, Luding and Behringer
(\emph{Phys.\ Rev.\ Lett.} 98, 2007) measured the contact numbers of photoelastic discs, which light up under
stress, and found a jump at $\varphi_J$ followed by growth with exponent about $0.5$, and pressure growing with
exponent about $1.1$ \PHYStag{experiment}. The square root, rather than the linear growth one might expect, says
that new contacts are cheap: the isostatic packing is marginal.

There is a pretty counting argument behind the marginality, due to Wyart, Nagel and Witten
(\emph{Europhys.\ Lett.} 72, 2005) \PHYStag{model}. Cut a cube of side $L$ (in units of the ball diameter)
out of a packing with $z=2d+\Delta z$ and free its boundary. The contacts crossing the boundary, about
$L^{d-1}$ of them, are lost, and in an isostatic packing each lost contact releases one floppy mode. The
interior has about $\frac12\Delta z\,L^d$ redundant contacts to block them. So floppy modes exist in boxes
of size up to
\[
  \ell^*\sim\frac1{\Delta z}\sim(\varphi-\varphi_J)^{-1/2},
\]
a length that diverges at the jamming point. Extended soft modes of this size explain the excess of
low-frequency vibrations (the ``boson peak'') of jammed packings.

\begin{figure}[ht]
\centering
\begin{tikzpicture}[x=55cm,y=0.48cm]
  \draw[->] (0.595,0)--(0.705,0) node[below right,font=\small]{$\varphi$};
  \draw[->] (0.6,0)--(0.6,10.5) node[left,font=\small]{$z$};
  \foreach \v in {4,6} {\draw (0.6,\v)--(0.598,\v) node[left,font=\scriptsize]{$\v$};}
  \draw (0.64,0)--(0.64,-0.25) node[below,font=\scriptsize]{$\varphi_J$};
  \draw[very thick,addcol] (0.6,0)--(0.64,0);
  \draw[very thick,addcol] plot[domain=0.64:0.70,samples=60] (\x,{6+14*sqrt(\x-0.64)});
  \fill[addcol] (0.64,6) circle (1.6pt);
  \draw[addcol,dashed] (0.64,0)--(0.64,6);
  \draw[thick,flagcol,dashed] (0.64,0)--(0.70,6.6) node[right,font=\scriptsize,flagcol]{$P\propto\varphi-\varphi_J$};
  \node[font=\scriptsize,addcol,anchor=west] at (0.70,9.45) {$z-6\propto(\varphi-\varphi_J)^{1/2}$};
  \draw[gray,densely dotted] (0.6,4)--(0.635,4);
  \node[font=\scriptsize,gray,anchor=west] at (0.6015,4.55) {friction: $z\ge4$};
  \node[font=\scriptsize,anchor=west] at (0.603,8.8) {$d=3$, schematic};
\end{tikzpicture}
\caption{Schematic behaviour of soft frictionless spheres near the jamming point in $d=3$: below $\varphi_J$
there are no contacts (after rattlers are removed); at $\varphi_J$ the mean contact number jumps to the
isostatic value $2d$; above it the excess grows like a square root, while the pressure of harmonic spheres
grows linearly. Frictional packings can be stable with fewer contacts.}
\label{PHYS:fig:jam}
\end{figure}

\subsection*{Friction, and the loosest random packing}
Real grains have friction, and friction lowers the number of contacts needed. The counting is the same as
in Corollary~\ref{SM:cor:maxwell}, but with forces instead of motions \PHYStag{model}. A frictional contact
carries a force \emph{vector} with $d$ components (normal and tangential), so a packing with mean
coordination $z$ has $\frac{zN}2\cdot d$ unknown force components. Each ball must satisfy force balance
($d$ equations) and torque balance ($\binom d2$ equations), $N\frac{d(d+1)}2$ in all. For generic packings
to carry a load we need at least as many unknowns as equations, $z\ge d+1$. Without friction a contact
carries only a normal force and there is no torque on a sphere, so $\frac{zN}2\ge dN$, which is $z\ge2d$
again. On the other hand each contact is one equation on the positions, whether or not there is friction,
so for generic sizes $\frac{zN}2\le dN$, as in Interlude~E. Hence
\[
  d+1\ \le\ z\ \le\ 2d,
\]
with $z=d+1$ approached for very rough grains and $z=2d$ for frictionless ones: $4\le z\le6$ in three
dimensions.

Fewer contacts allow looser packings. Onoda and Liniger (\emph{Phys.\ Rev.\ Lett.} 64, 1990) let glass
spheres settle slowly through liquids of nearly the same density and extrapolated to zero effective
gravity, finding a \emph{random loose packing} fraction $0.555\pm0.005$ \PHYStag{experiment}. Song, Wang and Makse
(\emph{Nature} 453, 2008) proposed a mean-field theory of jammed packings in which the packing fraction is
tied to the mean geometric coordination by $\varphi=z/(z+2\sqrt3)$ \PHYStag{model}. It gives
$\varphi=6/(6+2\sqrt3)\approx0.634$ for $z=6$ (frictionless, random close packing) and
$4/(4+2\sqrt3)\approx0.536$ for $z=4$ (infinite friction, random loose packing), and a phase diagram of
jammed matter between these two lines. With adhesion the constraint count drops further, and simulations of
small sticky grains give much looser stable packings; Liu, Li, Baule and Makse (\emph{Soft Matter} 11, 2015)
conjecture a limiting loosest packing at $z=2$ with $\varphi=1/2^3$, which their simulations approach but
do not reach \PHYStag{model}.

\subsection*{How packings are made}
Slow sedimentation gives loose packings; vibration, shearing and tapping compact them. Knight, Fandrich, Lau, Jaeger and Nagel (\emph{Phys.\ Rev.~E}
51, 1995) tapped glass beads in a tall narrow tube and found that the density approaches its final value
only logarithmically in the number of taps $t$ \PHYStag{experiment},
\[
  \varphi(t)=\varphi_\infty-\frac{\Delta\varphi_\infty}{1+B\ln(1+t/\tau)},
\]
with $\varphi_\infty$ between about $0.637$ and $0.647$ in the fits of Nowak, Knight, Ben-Naim, Jaeger and Nagel
(\emph{Phys.\ Rev.~E} 57, 1998). So tapping creeps towards random close packing and essentially stops there. Once the beads are annealed,
cycling the tapping intensity up and down retraces a reversible curve of density against intensity
(Nowak et al.), a standard testing ground for Edwards' statistical mechanics of powders
(Section~\ref{SM:sec:edwards}).

%======================================================================
\section{Dice, candies and Ulam's conjecture}
\label{PHYS:sec:shapes}

\subsection*{Counting for other shapes}
Lecture~6 ended with an aside on packing ellipsoids and polygons (Section~\ref{L04b:sec:aside}). In the
laboratory, non-spherical particles are the rule. The counting of Section~\ref{PHYS:sec:jam} goes through
for frictionless convex particles with $f$ degrees of freedom each: an isostatic packing has $2f$
constraints per particle on average. Spheres have $f=3$; spheroids (ellipsoids of revolution) $f=5$, since
rotation about the axis does nothing; general ellipsoids and tetrahedra $f=6$.

Donev, Cisse, Sachs, Variano, Stillinger, Connelly, Torquato and Chaikin (\emph{Science} 303, 2004) poured
M\&M candies, which are oblate spheroids, into a flask and measured a random packing fraction of
$0.685\pm0.01$, well above the $0.64$ of spheres \PHYStag{experiment}. Their simulations gave about $0.71$ for
spheroids and up to $0.735$ for general ellipsoids with aspect ratio about $1.3$ \PHYStag{simulation}.
The packing fraction has a cusp-like minimum at the sphere: any small deformation helps. The M\&Ms had
$9.82$ contacts on average, close to $2f=10$, but jammed general ellipsoids in the simulations had at most
about $11.4$, short of $2f=12$. Such \emph{hypostatic} packings are possible because the equivalence of jamming and first-order jamming,
which Interlude~E noted is special to spheres, fails for ellipsoids: the curvature of the surfaces can
block motions at second order, so a jammed packing of ellipsoids need not satisfy the first-order count.

\subsection*{Tetrahedral dice}
Aristotle believed that regular tetrahedra fill space, and Section~\ref{TET:sec} explained why they do not
even fill the space around an edge. Jaoshvili, Esakia, Porrati and Chaikin (\emph{Phys.\ Rev.\ Lett.} 104,
2010) poured tetrahedral dice, of the kind used in role-playing games, into containers and
measured, by volume and by magnetic resonance imaging of the interior, a random packing fraction of
$0.76\pm0.02$ \PHYStag{experiment}: denser than random spheres (by far) and even denser than the densest
sphere packing. A die touches only about $6.3$ neighbours on average, far below $2f=12$. But contacts between
polyhedra come in kinds: a face-to-face contact blocks one translation and two rotations, so counts as three
constraints, an edge-to-face contact as two, and a vertex-to-face or edge-to-edge contact as one. With these
weights the dice had about $12$ constraints per particle, isostatic again.

The densest packings of regular tetrahedra were a race between 2006 and 2010 \PHYStag{simulation}. The densest
lattice packing, of density $\frac{18}{49}\approx0.367$ (Hoylman 1970), is far from optimal; Conway and
Torquato (2006) found $0.72$, Chen (2008) $0.7786$, and then Haji-Akbari, Engel, Keys, Zheng, Petschek, Palffy-Muhoray and
Glotzer (\emph{Nature} 462, 2009), compressing tetrahedra on the computer, found that they
spontaneously form a \emph{dodecagonal quasicrystal}, a cousin of the cut-and-project sets of Interlude~C
(Section~\ref{ERG:sec:cutproject}), whose periodic approximant reaches $0.8503$. Within months the record
passed to explicit periodic packings built from pairs of tetrahedra glued along a face: $\frac{100}{117}\approx0.8547$
(Kallus, Elser and Gravel), $\frac{12250}{14319}\approx0.8555$ (Torquato and Jiao) and
$\frac{4000}{4671}\approx0.8563$ (Chen, Engel and Glotzer, gallery), the densest known. These constructions are exact,
but nobody knows whether $0.8563$ is optimal; the best upper bound is still astronomically close to $1$
(Gravel, Elser and Kallus, gallery). Lagarias and Zong (\emph{Notices AMS} 59, 2012) tell the story.

\subsection*{Is the ball the worst?}
Every one of these shapes packs more densely than the ball. Ulam suspected that this is always so: among
convex bodies in $\R^3$, the ball has the smallest optimal packing density (reported by Martin Gardner,
1995). This is \emph{Ulam's packing conjecture}, and it is open. Numerical work has ruled out many
candidates, including the regular tetrahedron. The best result is local: Kallus (gallery,
Theorem~\ref{PHYS:thm:kallus}) proved that every origin-symmetric convex body sufficiently close to the
ball, but not a ball, packs more densely than the ball, so the ball is a \emph{local pessimum}. In the plane
the disc is not the worst: the regular heptagon is conjectured to be, and among centrally symmetric convex
bodies Reinhardt (1934) conjectured that the worst is a smoothed octagon, of density
$\frac{8-4\sqrt2-\ln2}{2\sqrt2-1}\approx0.9024<\frac{\pi}{\sqrt{12}}\approx0.9069$. For centrally symmetric
bodies in the plane the densest packing is a lattice packing (the aside of Lecture~6), so Reinhardt's
problem is a problem about lattices; it is open too.

\begin{figure}[ht]
\centering
\begin{tikzpicture}
\begin{scope}[x=25cm,xshift=-8.25cm]
  \draw[->] (0.33,0)--(0.885,0) node[right,font=\small]{$\varphi$};
  \foreach \t in {0.35,0.40,0.45,0.50,0.55,0.60,0.65,0.70,0.75,0.80,0.85}
    {\draw (\t,0)--(\t,-0.08) node[below,font=\scriptsize]{$\t$};}
  % spheres row
  \fill[gray!15] (0.33,1.25) rectangle (0.494,1.55);
  \fill[addcol!35] (0.494,1.25) rectangle (0.545,1.55);
  \fill[addcol!15] (0.545,1.25) rectangle (0.7405,1.55);
  \node[font=\scriptsize] at (0.45,1.4) {fluid};
  \node[font=\scriptsize] at (0.68,1.4) {crystal};
  \node[font=\scriptsize,anchor=east] at (0.33,1.4) {spheres};
  \PHYSmsim{(0.384,1.0)} \node[font=\scriptsize,below] at (0.384,0.92) {RSA};
  \PHYSmsim{(0.494,1.0)} \PHYSmsim{(0.545,1.0)}
  \node[font=\scriptsize,above] at (0.5195,1.55) {coexistence};
  \PHYSmexp{(0.555,0.62)} \node[font=\scriptsize,below] at (0.555,0.55) {RLP};
  \PHYSmexp{(0.58,1.0)} \node[font=\scriptsize,below] at (0.58,0.92) {glass};
  \PHYSmexp{(0.64,1.0)} \PHYSmsim{(0.64,0.62)} \node[font=\scriptsize,below] at (0.64,0.55) {RCP, J};
  \PHYSmth{(0.7405,1.0)} \node[font=\scriptsize,below] at (0.7405,0.92) {FCC};
  % non-spherical row
  \node[font=\scriptsize,anchor=east] at (0.33,2.35) {other};
  \draw[gray!50] (0.33,2.35)--(0.885,2.35);
  \PHYSmcon{(0.340,2.35)} \node[font=\scriptsize,above] at (0.340,2.45) {diamond};
  \PHYSmth{(0.3673,2.35)} \node[font=\scriptsize,below] at (0.3673,2.25) {tetra.\ lattice};
  \PHYSmexp{(0.685,2.35)} \node[font=\scriptsize,above] at (0.685,2.45) {M\&M};
  \PHYSmsim{(0.735,2.35)} \node[font=\scriptsize,below] at (0.735,2.25) {ellipsoids};
  \PHYSmexp{(0.76,2.35)} \node[font=\scriptsize,above] at (0.765,2.45) {dice};
  \PHYSmsim{(0.8503,2.35)} \PHYSmcon{(0.8563,2.35)} \node[font=\scriptsize,above] at (0.853,2.45) {tetrahedra};
  % legend
  \PHYSmth{(0.36,3.2)} \node[font=\scriptsize,anchor=west] at (0.365,3.2) {theorem};
  \PHYSmcon{(0.46,3.2)} \node[font=\scriptsize,anchor=west] at (0.465,3.2) {exact construction};
  \PHYSmsim{(0.60,3.2)} \node[font=\scriptsize,anchor=west] at (0.605,3.2) {simulation};
  \PHYSmexp{(0.72,3.2)} \node[font=\scriptsize,anchor=west] at (0.725,3.2) {experiment};
\end{scope}
\end{tikzpicture}
\caption{Packing fractions of this interlude and of Interlude~E, by status. Spheres: random sequential
adsorption ($0.384$), the fluid--crystal coexistence region ($0.494$--$0.545$), random loose packing ($0.555$), the
colloidal glass (about $0.58$), random close packing and point~J (about $0.64$), and the Kepler bound
$\pi/\sqrt{18}$. Other shapes: the diamond packing of balls ($\pi\sqrt3/16$), the densest lattice packing of
regular tetrahedra ($18/49$), M\&M candies, simulated random ellipsoids, tetrahedral dice, the tetrahedral
quasicrystal approximant ($0.8503$) and the densest known tetrahedron packing ($4000/4671$).}
\label{PHYS:fig:strip}
\end{figure}

%======================================================================
\section{Gallery of theorems}
\label{PHYS:sec:gallery}

Results quoted in Interludes~E and~F (Hales, Hales--McLaughlin, Connelly et al., Richthammer, Theil) are not
repeated.

\begin{theorem}[J.~E.~Taylor, Ann.\ of Math.\ 103 (1976)]\label{PHYS:thm:taylor}
Every soap-film-like surface in $\R^3$, in the sense of Almgren's $(\mathbf M,\varepsilon,\delta)$-minimal sets, is
locally either a smooth minimal surface, or three smooth sheets meeting at $120^\circ$ along a curve, or
four such curves meeting at a point at angles $\arccos(-\frac13)$, with six sheets, in the shape of the cone over
the edges of a regular tetrahedron.
\end{theorem}
\noindent Plateau's experimental laws (Section~\ref{PHYS:sec:squash}) as a theorem. The tetrahedral angle is the
Tammes angle for $N=4$.

\begin{theorem}[T.~C.~Hales, Discrete Comput.\ Geom.\ 25 (2001)]\label{PHYS:thm:honeycomb}
Any partition of the plane into regions of equal area has perimeter at least that of the regular hexagonal
honeycomb (in the sense of an averaged limit over large discs).
\end{theorem}
\noindent The two-dimensional Kelvin problem is solved; the three-dimensional one is open.

\begin{question}[Kelvin's problem, 1887]
Which partition of $\R^3$ into cells of unit volume has the least area per cell? Is it the Weaire--Phelan
foam?
\end{question}
\noindent Kelvin's own candidate was beaten by $0.3\%$ in 1994, numerically but with a margin far beyond the
numerical error. That Weaire--Phelan is optimal is open.

\begin{theorem}[H.~Harborth, Elem.\ Math.\ 29 (1974)]\label{PHYS:thm:harborth}
The maximal number of touching pairs among $n$ non-overlapping unit discs in the plane is
$\lfloor3n-\sqrt{12n-3}\rfloor$.
\end{theorem}
\noindent The ground states of $n$ sticky discs. For $n=7$ (a hexagon around a disc) the formula gives
$21-9=12$.

\begin{theorem}[K.~Bezdek and S.~Reid, J.\ Geom.\ 104 (2013)]\label{PHYS:thm:bezdek}
Every packing of $n>2$ unit balls in $\R^3$ has fewer than $6n-0.926\,n^{2/3}$ touching pairs.
\end{theorem}
\noindent A lower bound for the sticky-sphere energy, with a surface term of the right order $n^{2/3}$. Exact values
are known only for $n\le5$ (Section~\ref{PHYS:sec:sticky}).

\begin{theorem}[E.~R.~Chen, M.~Engel and S.~C.~Glotzer, Discrete Comput.\ Geom.\ 44 (2010)]\label{PHYS:thm:CEG}
There is a periodic packing of regular tetrahedra of density $\frac{4000}{4671}\approx0.856347$, with four
tetrahedra (two face-sharing dimers) per fundamental domain. It is optimal within a three-parameter family of
dimer packings.
\end{theorem}
\noindent Found by numerical compression, then made exact; its global optimality is open. (Lattice packings,
whose tetrahedra all point the same way, reach only $\frac{18}{49}$: Hoylman, \emph{Bull.\ Amer.\ Math.\ Soc.}
76, 1970.)

\begin{theorem}[S.~Gravel, V.~Elser and Y.~Kallus, Discrete Comput.\ Geom.\ 46 (2011)]\label{PHYS:thm:GEK}
The packing density of regular tetrahedra is at most $1-2.6\cdots\times10^{-25}$.
\end{theorem}
\noindent A quantitative version of Section~\ref{TET:sec}; the gap between $0.856$ and $1-10^{-25}$ shows how
hard upper bounds are.

\begin{theorem}[Y.~Kallus, Adv.\ Math.\ 264 (2014)]\label{PHYS:thm:kallus}
There is $\varepsilon>0$ such that every origin-symmetric convex body $K\subset\R^3$ that is not a ball, and whose
asphericity is less than $\varepsilon$, has packing density greater than $\pi/\sqrt{18}$. In dimensions $4$ to $8$ and
$24$ there are origin-symmetric bodies of arbitrarily small asphericity whose lattice packing density is
smaller than that of the ball.
\end{theorem}
\noindent The best evidence for Ulam's conjecture. In dimensions $4$--$8$ and $24$ the best lattices for balls
($D_4,\dots,E_8$ of Lectures~4--6, and Leech) are so good that deforming the ball can make things worse.

\begin{theorem}[A.~Figotin and P.~Kuchment, SIAM J.\ Appl.\ Math.\ 56 (1996)]\label{PHYS:thm:FK}
Consider the periodic medium in $\R^d$, $d=2,3$, made of cubes of air, $[\delta,1]^d+\Z^d$, separated by walls of
width $\delta$ of a material with dielectric constant $\varepsilon>1$, and the operator $\nabla\cdot\varepsilon(x)^{-1}\nabla$. For any
$N$, if $\delta$ is small enough and $\varepsilon$ large enough (suitably coupled), the spectrum in $[0,N]$ lies in small
intervals around the Neumann eigenvalues of the cube. In particular the spectrum has gaps.
\end{theorem}
\noindent Band gaps for high contrast, as a theorem (for acoustic waves, or $2$-dimensional photonics); the
Maxwell case of inverse opals rests on computation (Section~\ref{PHYS:sec:opal}).

%======================================================================
\section*{Exercises}
\addcontentsline{toc}{section}{Exercises}

\begin{exercise}
(a) Show that the dihedral angle of the regular dodecahedron is $\arccos(-1/\sqrt5)$. (b) Show that the
truncated octahedron is the Voronoi cell of the body-centred cubic lattice $\Z^3\cup(\Z^3+\frac12(1,1,1))$, and
verify the values $5.345$ and $5.315$ of $A/V^{2/3}$ in Figure~\ref{PHYS:fig:rd} for it and for the rhombic
dodecahedron.
\end{exercise}

\begin{exercise}
(a) Use Proposition~\ref{PHYS:prop:coxeter} to find the average number of sides of a face for Kelvin's cell, for
the Weaire--Phelan cells, and for Matzke's bubbles. (b) The rhombic dodecahedron is not simple. Find the
analogue of the formula $F=12/(6-\bar n)$ for polyhedra in which every vertex has degree $3$ or $4$, and check it
on the rhombic dodecahedron. (c) Show that a convex polyhedron with three faces at each vertex and only
pentagonal and hexagonal faces has exactly $12$ pentagons.
\end{exercise}

\begin{exercise}
In the cell model of Proposition~\ref{PHYS:prop:cell}, at $\varphi=0.60$, compare the cage sizes of FCC and of an amorphous packing jamming at $0.64$, and
estimate the vibrational entropy difference per ball.
\end{exercise}

\begin{exercise}
(a) Show that the number of Barlow stackings of $L$ labelled layers is $3\cdot2^{L-1}$. (b) Show that in every
Barlow stacking each ball has $12$ touching neighbours, and that its Voronoi cell is either the rhombic or the
trapezo-rhombic dodecahedron. (c) Using Proposition~\ref{PHYS:prop:planes}, find the second largest spacing of
lattice planes in FCC, and the corresponding Bragg wavelength of an opal at normal incidence.
\end{exercise}

\begin{exercise}
(a) Redo Proposition~\ref{PHYS:prop:letokhov} for a cube of side $L$ (Dirichlet conditions on its boundary),
and show that the threshold is $L=\pi\sqrt{\ell\ell_g}$.
(b) With $\ell=1/(n\sigma)$ and $n=6\varphi/(\pi D^3)$, how does the critical number of grains in a ball of
radius $R_{\mathrm{cr}}$ depend on $\varphi$?
\end{exercise}

\begin{exercise}
(a) Check that the six points listed in Example~\ref{PHYS:ex:six} form three regular tetrahedra glued face to face,
and that they have exactly $12$ contacts. (b) Verify $\det I=8$ and $\det(AA^T)=864$ for the octahedron.
(\emph{Hint:} $AA^T=2\,\mathrm{Id}+B$, where $B$ records the angles between contacts sharing a ball.) (c) For $n=5$
show that a rigid cluster must be the triangular bipyramid, with $9$ contacts.
\end{exercise}

\begin{exercise}
(a) Carry out the frictional counting of Section~\ref{PHYS:sec:jam} for discs in the plane: show that
$3\le z\le4$. (b) For frictionless tetrahedra, explain why the
isostatic count is $12$ constraints per particle, and check it against the contact weights of
Section~\ref{PHYS:sec:shapes}.
\end{exercise}

\begin{exercise}
Check Harborth's formula for the hexagonal patches $n=3k^2-3k+1$: show that such a patch has $3(3k^2-5k+2)$
contacts and that this equals $\lfloor3n-\sqrt{12n-3}\rfloor$.
\end{exercise}

\section*{Further reading}
\addcontentsline{toc}{section}{Further reading}
\begin{itemize}[nosep]
\item H.~S.~M.~Coxeter, \emph{Close-packing and froth}, Illinois J.\ Math.\ 2 (1958). The statistical honeycombs of
  Proposition~\ref{PHYS:prop:coxeter}.
\item D.~Weaire (ed.), \emph{The Kelvin Problem: Foam Structures of Minimal Surface Area}, Taylor and Francis, 1996;
  D.~Weaire and S.~Hutzler, \emph{The Physics of Foams}, Oxford, 1999.
\item P.~N.~Pusey et al., \emph{Hard spheres: crystallization and glass formation}, Phil.\ Trans.\ R.\ Soc.\ A 367
  (2009), arXiv:0908.1663.
\item J.~D.~Joannopoulos, S.~G.~Johnson, J.~N.~Winn and R.~D.~Meade, \emph{Photonic Crystals: Molding the Flow of
  Light}, 2nd ed., Princeton, 2008.
\item D.~S.~Wiersma, \emph{The physics and applications of random lasers}, Nature Physics 4 (2008).
\item M.~Holmes-Cerfon, \emph{Sticky-sphere clusters}, Annu.\ Rev.\ Condens.\ Matter Phys.\ 8 (2017),
  arXiv:1709.05138; and \emph{Enumerating rigid sphere packings}, SIAM Review 58 (2016).
\item A.~J.~Liu and S.~R.~Nagel, \emph{Jamming is not just cool any more}, Nature 396 (1998).
\item J.~C.~Lagarias and C.~Zong, \emph{Mysteries in packing regular tetrahedra}, Notices AMS 59 (2012).
\end{itemize}
